\subsection{Riconduzione ai limiti notevoli}
\label{sec:esercizi-limiti-limiti-notevoli}

\begin{exercise}[Seno di un multiplo dell'argomento]
\label{ex:limiti-foglio-2311-18}
\difficoltaesercizio{0}
Calcolare il limite
\[
  \lim_{x\to0}\frac{\sin(\pi x)}{\pi x}.
\]
\rimandoteoria{thm:lez20-limite-seno}
\end{exercise}
\begin{proof}[Soluzione dell'esercizio]
Posto $t:=\pi x$, si ha $t\to0$ per $x\to0$.
Per il limite notevole del seno,
\[
  \lim_{x\to0}\frac{\sin(\pi x)}{\pi x}
  =\lim_{t\to0}\frac{\sin t}{t}=1.
\]
\end{proof}

\begin{exercise}[Seno della tangente]
\label{ex:limiti-foglio-2311-19}
\difficoltaesercizio{0}
Calcolare il limite
\[
  \lim_{x\to0}\frac{\sin(\tan x)}{\sin x}.
\]
\rimandoteoria{thm:lez20-limite-seno}
\end{exercise}
\begin{proof}[Soluzione dell'esercizio]
Poiché $\tan x\to0$ per $x\to0$, per il limite notevole
del seno e la continuità del coseno si ottiene
\[
\begin{aligned}
  \lim_{x\to0}\frac{\sin(\tan x)}{\sin x}
    &=\lim_{x\to0}\frac{\sin(\tan x)}{\tan x}\cdot\frac1{\cos x}\\
    &=1.
\end{aligned}
\]
\end{proof}

\begin{exercise}[Differenza di seni divisa per la variabile]
\label{ex:limiti-foglio-2311-20}
\difficoltaesercizio{0}
Calcolare il limite
\[
  \lim_{x\to0}\frac{\sin(5x)-\sin(3x)}{x}.
\]
\rimandoteoria{thm:lez20-limite-seno}
\end{exercise}
\begin{proof}[Soluzione dell'esercizio]
Riscrivendo i due addendi e applicando il limite notevole del seno,
\[
\begin{aligned}
  \lim_{x\to0}\frac{\sin(5x)-\sin(3x)}{x}
    &=\lim_{x\to0}\left(5\frac{\sin(5x)}{5x}
      -3\frac{\sin(3x)}{3x}\right)\\
    &=5-3=2.
\end{aligned}
\]
\end{proof}

\clearpage

\begin{exercise}[Coseno di un multiplo dell'argomento]
\label{ex:limiti-foglio-2311-24}
\difficoltaesercizio{0}
Calcolare il limite
\[
  \lim_{x\to0}\frac{1-\cos(2x)}{x^2}.
\]
\rimandoteoria{prop:lez20-limite-coseno}
\end{exercise}
\begin{proof}[Soluzione dell'esercizio]
Posto $t:=2x$, si ha $t\to0$ per $x\to0$.
Per il limite notevole del coseno,
\[
  \lim_{x\to0}\frac{1-\cos(2x)}{x^2}
  =4\lim_{t\to0}\frac{1-\cos t}{t^2}
  =4\cdot\frac12=2.
\]
\end{proof}
